% from ega4-15-fr.tex
\oldpage[IV-3]{227}
\begin{lemma}[15.2.2.1]
\label{IV.15.2.2.1-fr}
Soient $Y=\operatorname{Spec}(A)$ le spectre d'un anneau de valuation discrète,
$y$ son point fermé, $f:X\to Y$ un morphisme localement de type fini, $x$ un
point de $f^{-1}(y)$, $\mathcal F$ un $\mathcal O_X$-Module cohérent. On suppose
vérifiées les conditions suivantes~:
\begin{enumerate}[label=(\roman*)]
  \item $\operatorname{Supp}(\mathcal F)=X$.
  \item $f$ est ouvert aux points génériques des composantes irréductibles de
  $f^{-1}(y)$ contenant $x$.
  \item $\mathcal F_y=\mathcal F\otimes_{\mathcal O_y}\mathbf{k}(y)$ est réduit
  au point $x$ (3.2.2).
\end{enumerate}
Alors $\mathcal F$ est $f$-plat au point $x$.
\end{lemma}

Désignons en effet par $t$ une uniformisante de $A$, posons $B=\mathcal O_x$ et
$\mathcal F_x=M$~; la condition (i) implique que
$\operatorname{Supp}(M)=\operatorname{Spec}(B)$ (\textbf{I}, 9.1.13). La
condition (ii) implique, en vertu de (14.3.2), que toute composante irréductible
de $X$ contenant $x$ domine $Y$~; en d'autres termes, l'image réciproque dans $A$
de tout idéal premier minimal $\mathfrak p_i$ de $B$ est $0$, ce qui signifie
encore que l'on a $t\notin\mathfrak p_i$ pour tout $i$. Enfin, (iii) signifie que
$M/tM$ est un $(B/tB)$-module \emph{réduit} (3.2.2). Il s'agit de montrer que
sous ces conditions $M$ est un $A$-module \emph{sans torsion}
($\mathbf{0}_{\mathrm I}$, 6.3.4), et pour cela il suffit évidemment de montrer
que $t$ est un élément $M$-\emph{régulier}~; mais cela résulte de (3.4.6.1).
\end{proof}

\begin{corollary}[15.2.3]
\label{IV.15.2.3-fr}
Soient $Y$ un préschéma localement noethérien, $f:X\to Y$ un morphisme localement
de type fini, $x$ un point de $X$, $y=f(x)$. On suppose vérifiées les conditions
suivantes~:
\begin{enumerate}[label=(\roman*)]
  \item $f$ est universellement ouvert aux points génériques des composantes
  irréductibles de $f^{-1}(y)$ contenant $x$.
  \item $f^{-1}(y)$ est géométriquement réduit (sur $\mathbf{k}(y)$) au point
  $x$ (4.6.9).
  \item L'anneau $\mathcal O_y$ est réduit.
\end{enumerate}
Sous ces conditions $f$ est plat au point $x$.
\end{corollary}

Il suffit d'appliquer (15.2.2) à $\mathcal F=\mathcal O_X$ (4.6.22).

\begin{remarks}[15.2.4]
\label{IV.15.2.4-fr}
\begin{enumerate}[label=(\roman*)]
  \item Les trois premières des conditions de (15.2.2) ne changent pas lorsqu'on
  remplace $f$ par $f_{\mathrm{red}}$~; mais si $\mathfrak N$ est le nilradical
  d'un anneau local noethérien $A$, $A/\mathfrak N$ est plat sur $A/\mathfrak N$
  mais non sur $A$ lui-même lorsque $\mathfrak N\ne0$ (Bourbaki,
  \emph{Alg. comm.}, chap.~II, \S~3, n\textsuperscript{o}~2, cor.~2 de la
  prop.~5). On voit donc qu'on ne peut dans (15.2.2) supprimer la condition (iv).

  \item Il résulte de (2.4.6) que si la conclusion de (15.2.2) est vraie, ainsi
  que l'hypothèse (i), $f$ est universellement ouvert dans un voisinage de $x$,
  et en particulier aux points génériques des composantes irréductibles de
  $f^{-1}(y)$ contenant $x$. Même lorsque (i), (iii) et (iv) sont vérifiées, on
  ne peut remplacer (ii) par l'hypothèse plus faible que $f$ est universellement
  ouvert \emph{au point $x$ seulement}. C'est ce que montre l'exemple
  (14.1.3, (i)), où $f$ est évidemment universellement ouvert en tout point de
  $X_2$, donc en particulier au point $x$, intersection de $X_1$ et $X_2$~; les
  conditions (ii) et (iii) de (15.2.3) sont aussi vérifiées par cet exemple.
  Toutefois, $\mathcal O_x$ n'est pas un $\mathcal O_y$-module sans torsion,
  $\mathcal O_y$ s'identifiant à un sous-anneau de $\mathcal O_x$ qui contient
  des diviseurs de $0$ dans $\mathcal O_x$~; donc $f$ n'est pas plat au point $x$.

  \item Dans (15.2.3), la conclusion n'est plus nécessairement valable lorsqu'on
  remplace l'hypothèse (ii) par l'hypothèse plus faible que $f^{-1}(y)$,
  considéré comme $\mathbf{k}(y)$-préschéma, n'a pas de cycles premiers associés
  immergés. Un exemple est fourni en prenant pour $Y$ une courbe ayant un point
  de rebroussement $y$, par exemple
  $Y=\operatorname{Spec}(K[S,T]/(S^3-T^2))$, $K$ étant un corps algébriquement
  clos, et pour $X$ sa normalisée (\textbf{II}, 6.3.8). Si $s$ et $t$ sont les
  classes de $S$ et $T$ dans $A=K[S,T]/(S^3-T^2)$, on vérifie aussitôt que
  $X=\operatorname{Spec}(B)$, où $B=K[s,t,u]$, $u$ étant l'élément $t/s$ du corps
  des fractions de $A$, et comme $s=u^2$, $t=u^3$, on a aussi $B=K[u]$,
  isomorphe à l'anneau de polynômes à une indéterminée sur $K$. Le seul point $x$
  de $X$ au-dessus du point $y$ correspondant à l'idéal maximal $(s)+(t)$
  correspond à l'idéal maximal $(u)$~; mais comme la classe $\bar u$ de $u$ dans
  $\mathcal O_x/\mathfrak m_y\mathcal O_x$ est telle que $\bar u^2=0$,
  $f^{-1}(y)$ n'est pas un $\mathbf{k}(y)$-préschéma réduit. Ici $f$ est un
  morphisme fini, surjectif et radiciel, donc un homéomorphisme universel (2.4.5),
  mais n'est pas plat au point $x$, car si $\mathcal O_x$ était un
  $\mathcal O_y$-module plat, il serait un $\mathcal O_y$-module
  